\documentclass[12pt,reqno]{amsart}
\usepackage{math-review}
% Project-wide mathematical notation. No claims are certified here.

\title[Square completion and the Schur complement]{Deriving the Schur-complement criterion\\by square completion}
\author{A lecture on classical linear algebra}
\date{}
\begin{document}
\maketitle

\section{The goal and the operation that will prove it}
A block quadratic form couples two vectors through a mixed term. The aim
is to remove that coupling by a change of variables and then test the two
remaining terms separately. The complete goal is the following classical
criterion.
\begin{theorem}\label{thm:main}
Let $p,q\geq1$ be integers. Let $A\in\mathbb R^{p\times p}$ and
$C\in\mathbb R^{q\times q}$ be symmetric, and let $B\in\mathbb R^{p\times q}$.
Assume that $x^{\mathsf T}Ax>0$ for every $0\ne x\in\mathbb R^p$. Define
\[
 M=\begin{pmatrix}A&B\\ B^{\mathsf T}&C\end{pmatrix},
 \qquad S=C-B^{\mathsf T}A^{-1}B.
\]
Then $M$ is positive definite if and only if $S$ is positive definite.
Here a real symmetric matrix $H$ is positive definite when
$v^{\mathsf T}Hv>0$ for every nonzero real vector $v$ of the appropriate size.
\end{theorem}
Throughout, $H>0$ abbreviates the positive definiteness defined in the
statement. The fixed data are the blocks $A,B,C$; the vectors $x,y$ will
range freely over their respective real spaces. Multiplying the blocks gives
\begin{equation}\label{eq:Q}
 Q(x,y):=\begin{pmatrix}x\\y\end{pmatrix}^{\!\mathsf T}
 M\begin{pmatrix}x\\y\end{pmatrix}
 =x^{\mathsf T}Ax+2x^{\mathsf T}By+y^{\mathsf T}Cy.
\end{equation}
The cross term $2x^{\mathsf T}By$ prevents a direct separate sign test.
We will choose $z=x+Ky$ to absorb it into $z^{\mathsf T}Az$. Matching the
cross term will force $AK=B$ and hence $K=A^{-1}B$.

There are two tasks after this choice: compute the remainder exactly, and
verify that $(x,y)\mapsto(z,y)$ preserves nonzero pairs. The first is
Lemma~\ref{lem:square}; the second is Lemma~\ref{lem:inverse}. Together they
prove sufficiency. Necessity uses the particular choice $x=-A^{-1}By$,
which cancels the first term for each nonzero $y$.

\paragraph{Reader inputs.}
We assume multiplication of compatible real matrices, the transpose rule
$(UV)^{\mathsf T}=V^{\mathsf T}U^{\mathsf T}$, and expansion of a scalar
square. Positive definiteness was recalled in the theorem. We also import
the elementary finite-dimensional fact that a square linear map with zero
kernel is invertible. We will use it only to justify $A^{-1}$; the block
identity and its inverse change of variables are proved below.
No Schur-complement or congruence theorem is assumed.

\section{Choose the shift by matching the mixed term}
For real scalars $a>0,b,c$, expand a trial square:
\[
 a(x+ky)^2=ax^2+2akxy+ak^2y^2.
\]
Matching $2bxy$ means $ak=b$. Substituting $k=b/a$ leaves
\begin{equation}\label{eq:scalar}
 ax^2+2bxy+cy^2
 =a\left(x+\frac ba y\right)^2+
   \left(c-\frac{b^2}{a}\right)y^2.
\end{equation}
The last coefficient is the cost of removing the mixed term. In blocks
we seek a $p\times q$ matrix $K$, so that $Ky\in\mathbb R^p$ can be
added to $x$. Symmetry of $A$ gives
\[
 (x+Ky)^{\mathsf T}A(x+Ky)
 =x^{\mathsf T}Ax+2x^{\mathsf T}AKy+y^{\mathsf T}K^{\mathsf T}AKy.
\]
Thus $AK=B$ matches the mixed terms for all $x,y$, exactly as $ak=b$
did in the scalar calculation.

Here the hypothesis $A>0$ first enters. If $Au=0$, then
$u^{\mathsf T}Au=0$, so positivity forces $u=0$. The imported kernel test
now gives invertibility, and the unique choice is
\[
 K=A^{-1}B\in\mathbb R^{p\times q},\qquad z=x+A^{-1}By.
\]
The construction has supplied an admissible vector $z$ of the correct
size. The next calculation identifies exactly what remains.

\section{Expand the square and identify the remainder}
Symmetry and invertibility give
$(A^{-1})^{\mathsf T}=(A^{\mathsf T})^{-1}=A^{-1}$. Hence
\begin{align*}
 z^{\mathsf T}Az
 &=(x+A^{-1}By)^{\mathsf T}A(x+A^{-1}By)\\
 &=x^{\mathsf T}Ax+x^{\mathsf T}AA^{-1}By
   +y^{\mathsf T}B^{\mathsf T}(A^{-1})^{\mathsf T}Ax\\
 &\hspace{1.5em}+y^{\mathsf T}B^{\mathsf T}(A^{-1})^{\mathsf T}AA^{-1}By\\
 &=x^{\mathsf T}Ax+x^{\mathsf T}By+y^{\mathsf T}B^{\mathsf T}x
   +y^{\mathsf T}B^{\mathsf T}A^{-1}By\\
 &=x^{\mathsf T}Ax+2x^{\mathsf T}By+y^{\mathsf T}B^{\mathsf T}A^{-1}By.
\end{align*}
The middle terms coincide because the scalar $x^{\mathsf T}By$ equals
its transpose. Subtract the last term and add $y^{\mathsf T}Cy$.
This proves the exact identity we need, with weaker assumptions than the
goal theorem.
\begin{lemma}\label{lem:square}
Let $A\in\mathbb R^{p\times p}$ be symmetric and invertible,
$B\in\mathbb R^{p\times q}$ and $C\in\mathbb R^{q\times q}$ symmetric.
For all $x\in\mathbb R^p$ and $y\in\mathbb R^q$, putting
$z=x+A^{-1}By$ gives
\[
 x^{\mathsf T}Ax+2x^{\mathsf T}By+y^{\mathsf T}Cy
 =z^{\mathsf T}Az+y^{\mathsf T}(C-B^{\mathsf T}A^{-1}B)y.
\]
\end{lemma}
Thus the remaining matrix is $S=C-B^{\mathsf T}A^{-1}B$, the Schur
complement of $A$ in $M$. It is a symmetric $q\times q$ matrix:
\[
 S^{\mathsf T}=C^{\mathsf T}-B^{\mathsf T}(A^{-1})^{\mathsf T}B=S.
\]
The square calculation used symmetry and invertibility, without using
the sign of $A$. That sign becomes essential in the positivity argument.

\section{Preserve nonzero vectors and prove both directions}
An identity in new variables is useful only if the change does not lose
any vectors that positivity must test. The precise elementary fact is:
\begin{lemma}\label{lem:inverse}
For every real matrix $K\in\mathbb R^{p\times q}$, the linear map
\[
 (x,y)\longmapsto (x+Ky,y)
\]
is invertible, with inverse $(z,y)\mapsto(z-Ky,y)$. Consequently it takes
nonzero vectors to nonzero vectors, in both directions.
\end{lemma}
\begin{proof}
The two compositions are
\[
 (x,y)\longmapsto(x+Ky,y)\longmapsto((x+Ky)-Ky,y)=(x,y)
\]
and
\[
 (z,y)\longmapsto(z-Ky,y)\longmapsto((z-Ky)+Ky,y)=(z,y).
\]
Both maps take zero to zero. Applying the inverse shows that a nonzero
pair cannot map to zero. The same argument applies in the reverse direction.
\end{proof}
We use this lemma with $K=A^{-1}B$. The proof now separates into two
arguments with different choices of vectors.

\begin{proof}[Proof of Theorem~\ref{thm:main}]
\emph{Sufficiency.} Assume $S>0$, retaining the theorem's hypothesis $A>0$.
Take an arbitrary nonzero pair $(x,y)$ and set $z=x+A^{-1}By$.
Lemma~\ref{lem:inverse} ensures that $(z,y)\ne(0,0)$, and
Lemma~\ref{lem:square} gives
\[
 Q(x,y)=z^{\mathsf T}Az+y^{\mathsf T}Sy.
\]
If $z\ne0$, the first term is positive and the second is nonnegative.
If $z=0$, then $y\ne0$, so the second term is positive and the first
is zero. In either case $Q(x,y)>0$. Since the original nonzero pair was
arbitrary, this proves $M>0$.

\emph{Necessity.} Assume $M>0$ and let $y\in\mathbb R^q$ be an arbitrary
nonzero vector. Choose $x=-A^{-1}By$. This makes $z=0$, while $(x,y)$
is still nonzero because its second component is nonzero. Therefore
\[
 y^{\mathsf T}Sy=Q(-A^{-1}By,y)>0.
\]
This holds for every nonzero $y$, which is exactly $S>0$.
\end{proof}
The inverse map protects the universal test in the forward direction;
the cancellation vector isolates the residual form in the reverse
direction. No determinant criterion or eigenvalue formula was needed.

\section{Execute the method and locate its limit}
Take $p=q=1$, $A=2$, $B=1$, $C=1$. The prescribed operations give
\[
 A^{-1}B=\frac12,\qquad S=1-\frac12=\frac12,
\]
and the complete quadratic-form calculation is
\[
 2x^2+2xy+y^2=2(x+y/2)^2+\frac12y^2.
\]
For $(x,y)\ne(0,0)$, the two terms cannot both vanish: vanishing of the
second forces $y=0$, and vanishing of the first then forces $x=0$.
This verifies positive definiteness directly. With $y$ fixed, the first
term has unique minimum zero at $x=-y/2$, leaving minimum $y^2/2$.

To see why invertibility alone is insufficient, take $A=-1$, $B=0$,
$C=1$. Then $A$ is invertible and $S=1>0$, but
\[
 Q(x,y)=-x^2+y^2,\qquad Q(1,0)=-1<0.
\]
The square identity remains valid. The precise failing step is the
assertion that $z^{\mathsf T}Az$ is positive when $z\ne0$.
This is a counterexample to dropping $A>0$, not to the stated theorem.
If $A$ is singular, $A^{-1}B$ is unavailable, so this construction does
not apply. No pseudoinverse criterion is asserted.

The exposition follows the elementary model in the supplied development
guide~\cite[Appendix A]{S-guide}, with the square expansion, inverse map,
and both positivity directions written out here. It concerns known
finite-dimensional linear algebra and makes no originality claim.
\bibliographystyle{plain}
\bibliography{../generated/references}
\end{document}
